\begin{proposition}[Open question and proved period identities]\label{ari:C.7.4}
\StatementRecord{ari:C.7.4}{open; current checks have the scopes recorded below}{Historical records: GPT-6 Astra (OpenAI, ChatGPT Pro mode), effort Pro, 2026-11-17 (historical check record); Claude Fable 5.2 (Anthropic), effort high, 2026-11-17. Also checked in this writing by \CurrentChecker. Re-derived the Legendre two-by-two inversion, real-linear quasi-period formula, compact-character equation, and nonzero central sum functional on plane and positive bracket line. Receiving text remains open for transcendence and states only these proved identities; the former conjecture-dependent argument is retained only in the unchanged compendium.}{\cref{ari:C.7.2,ari:C.7.3,ari:eq:cp-v,ari:eq:cp-positive-centre};
the Legendre relation and the displayed two-by-two linear systems;
\cite[Definitions 2.3--2.4 and Remark 4.13]{bertolin2022}}

For an algebraic parameter $t\ne0,1$ with nontorsion $P_t$, the
transcendence of $N(P_t)$ remains open in this writing. The following
identities are proved. Write $P=[u]$ on the algebraic elliptic curve
$E=\C/(\Z\omega_1+\Z\omega_2)$, with quasi-periods $\eta_1,\eta_2$,
and use the extension $G_P$, logarithmic coordinate $v$, and compact
character of \cref{ari:C.7.2}. The field $K_0$ generated over
$\overline\Q$ by the full de Rham--Betti period matrices of
$G_P\times G_{\bar P}$ contains $u,\zeta(u),\eta(u)$ and their
complex conjugates. Every compact lift has
\begin{equation}\label{ari:eq:cp-relation}
 v+\bar v=u\zeta(u)+\bar u\overline{\zeta(u)}
              -u\eta(u)-\bar u\overline{\eta(u)}.
\end{equation}
The linear functional $(c_1,c_2)\mapsto c_1+c_2$ is nonzero on
each central bracket space computed in \cref{ari:C.7.3}.
These are the established period identities recorded for the open
transcendence question.
\end{proposition}

\begin{proof}[Proof of the period identities]
\begin{enumerate}[label=\arabic*.,leftmargin=*]
\item The full period matrix contains the pure periods, quasi-periods,
$2\pi i$, and the third-kind periods
\[
 \kappa_j=\zeta(u)\omega_j-\eta_j u,\qquad j=1,2.
\]
The two-by-two coefficient matrix for $(\zeta(u),u)$ is
$\left(\begin{smallmatrix}\omega_1&-\eta_1\\
\omega_2&-\eta_2\end{smallmatrix}\right)$.
Its determinant is the nonzero Legendre determinant, up to sign,
$2\pi i$. Inverting this matrix proves that $u$ and $\zeta(u)$
belong to $K_0$. Applying the same computation to the conjugate
extension places their conjugates in that field as well.

\item Write $u=a\omega_1+b\omega_2$ with $a,b\in\R$. Solving this
system together with its complex conjugate gives
\begin{equation}\label{ari:eq:cp-real-linear}
 \eta(u)=
 \frac{(u\bar\omega_2-\bar u\omega_2)\eta_1
       +(\bar u\omega_1-u\bar\omega_1)\eta_2}
      {\omega_1\bar\omega_2-\bar\omega_1\omega_2}.
\end{equation}
The denominator is nonzero because $\omega_1,\omega_2$ are
real-linearly independent. This is the actual real-linear extension
of $\omega_j\mapsto\eta_j$, and it proves the claimed field
membership without identifying complex conjugation with a
rational function of the unconjugated periods.

\item The compact-character formula of \cref{ari:C.7.2} is
$\operatorname{Re}(v-u\zeta(u)+u\eta(u))=0$.
Adding this equation to its conjugate yields
\eqref{ari:eq:cp-relation}. The right side belongs to $K_0$ by
the preceding two computations.

\item For the two-dimensional bracket space of \cref{ari:C.7.3},
the sum functional takes value one on $(1,0)$.
For its one-dimensional bracket space, the specified generator
is $(1,c)$, where $c=aa^\dagger/d>0$ by the actual polarization
transport through the degree-$d$ isogeny. Its value is $1+c>0$.
Any additional central Tate line leaves this nonzero restriction
unchanged. Thus the displayed equation cuts a proper affine
hyperplane in each of these explicit central coordinate spaces.

\item The question asks whether the local norm at the stated algebraic
points is transcendental. The identities just proved give no
unconditional transcendence-degree lower bound. The earlier
argument invoking an unresolved period conjecture is preserved,
with its premise, in the unchanged S6\_09 proof compendium under
the same label; it is not used as a theorem of this paper.
\end{enumerate}
\end{proof}